\chapter{پاداش، یادگیری تقویتی و تصمیم‌گیری}

\section{سیستم پاداش و یادگیری تقویتی}
پس از پایان بحث مدل‌های شنیداری، بحث وارد موضوع سیستم‌های پاداش در مغز و ارتباط آن با \eng{Reinforcement Learning} شد. نقطهٔ شروع این بحث این است که یک عامل در محیط قرار دارد، در هر لحظه در یک وضعیت یا \eng{state} است، عملی یا \eng{action} انجام می‌دهد و از یک وضعیت به وضعیت دیگر منتقل می‌شود. به این تغییر وضعیت \eng{transition} گفته می‌شود.

هر \eng{transition} می‌تواند با پاداش یا تنبیه همراه باشد. پاداش ممکن است مثبت، منفی، کم، زیاد یا خنثی باشد. این پاداش‌ها نقش \eng{reinforcer} دارند: باعث می‌شوند عامل یاد بگیرد کدام رفتارها را بیشتر تکرار کند و کدام رفتارها را کم‌تر انجام دهد.

هدف این بخش ورود سنگین به معادلات فیزیکی یا ریاضی نبود. ابتدا قرار بود سازوکار کلی مغز در یادگیری از طریق بازخورد فهمیده شود و سپس این بحث به فضای \eng{Reinforcement Learning} وصل شود.

در معماری شناختی‌ای که پیش‌تر دربارهٔ آن صحبت شده بود، عامل از محیط محرک‌های حسی دریافت می‌کند. این محرک‌ها، برای مثال از مسیر ماژول بینایی یا شنیداری، وارد سیستم می‌شوند. سپس با فرایندهایی مانند توجه انتخابی، بخشی از اطلاعات محیط انتخاب می‌شود. اطلاعات فعلی وارد حافظهٔ کوتاه‌مدت می‌شود؛ سیستم از حافظهٔ بلندمدت اطلاعات مرتبط را بازیابی می‌کند؛ با ترکیب اطلاعات فعلی و دانش قبلی، یک الگو استخراج می‌شود؛ و بر اساس آن، عمل انتخاب و از مسیر ماژول حرکتی اجرا می‌شود.

پس از انجام عمل، عامل از زمان \(t\) به زمان \(t+1\) می‌رود و به وضعیت جدیدی منتقل می‌شود. در همین فاصله، پاداش یا تنبیه دریافت می‌شود. در این نقطه باید چیزی شبیه یک \eng{utility module} یا ماژول ارزش‌گذاری به معماری اضافه شود؛ ماژولی که بررسی کند عمل انجام‌شده خوب بود، بد بود، یا تفاوت خاصی با انتظار قبلی نداشت. هدف این ماژول این است که پس از هر تجربه، مدل قبلی یا \eng{prior model} به‌روزرسانی شود.

\subsection{پاداش در مغز}
در انسان، پاداش معمولاً به‌صورت حس رضایت، لذت یا موفقیت تجربه می‌شود. غذا خوردن، موفق شدن در یک کار، دریافت پیام مهم یا دیدن چیزی که قبلاً تجربهٔ خوبی با آن داشته‌ایم، نمونه‌هایی از تجربه‌های پاداش‌دهنده‌اند.

این تجربه‌ها می‌توانند با آزاد شدن دوپامین همراه شوند. نکتهٔ مطرح‌شده در این بحث این بود که دوپامین در واکنش به یک تغییر حالت معنادار آزاد می‌شود؛ تغییری که برای سیستم مهم یا آموخته‌شده است. برای مثال خاموش شدن چراغ در روز ممکن است تجربهٔ خاصی نباشد، اما دریافت یک پیامک به‌دلیل انتظار یا تجربهٔ آموخته‌شده می‌تواند برای فرد پاداش‌دهنده باشد.

\subsection{دوپامین و \eng{VTA}}
وقتی دوپامین آزاد می‌شود، در مغز از مسیرهایی عبور می‌کند که به آن‌ها مسیرهای دوپامینرژیک یا \eng{dopaminergic pathways} گفته می‌شود. در بحث پاداش، یکی از نقاط مهم تولید دوپامین ناحیهٔ \eng{Ventral Tegmental Area} یا \eng{VTA} است.

در چارچوب این بخش، کارکرد مهم \eng{VTA} محاسبهٔ خطای پیش‌بینی پاداش است: سیستم چه پاداشی را پیش‌بینی کرده بود، در واقعیت چه پاداشی دریافت شد، و اختلاف این دو چقدر است. وقتی پاداش واقعی بهتر از انتظار باشد، \eng{VTA} با ترشح دوپامین مرتبط می‌شود. بنابراین دوپامین صرفاً نشانهٔ «لذت» نیست؛ بلکه در این بحث بیشتر به غافلگیری مثبت و خطای پیش‌بینی پاداش مربوط است.

به همین دلیل، هر چیز خوب یا لذت‌بخشی الزاماً همیشه دوپامین زیادی ایجاد نمی‌کند. اگر تجربه‌ای کاملاً قابل پیش‌بینی شده باشد، ممکن است غافلگیری کمی ایجاد کند و در نتیجه پاسخ دوپامینی آن نیز کم‌تر شود.

\subsection{مسیرهای دوپامینرژیک}
چهار مسیر اصلی دوپامینرژیک در مغز مطرح می‌شود، اما در بحث پاداش این بخش تمرکز بر دو مسیر بود:
\[
\text{\lr{VTA}}\rightarrow
\begin{cases}
\text{\lr{Mesolimbic Pathway}}\\
\text{\lr{Mesocortical Pathway}}
\end{cases}
\]

\coursefigure{0.62}{dopaminergic-vta-pathways.png}{مسیرهای دوپامینرژیک مرتبط با \eng{VTA}: ارتباط ناحیهٔ تگمنتوم شکمی با هستهٔ اکومبنس، استریاتوم، قشر پیشانی، هیپوکمپ و تودهٔ سیاه.}

\subsubsection{\eng{Mesolimbic Pathway}}
مسیر \eng{Mesolimbic} از \eng{VTA} به سمت سیستم لیمبیک می‌رود. این مسیر با احساس رضایت، شرطی‌شدن، تقویت رفتار، اعتیاد و میل به تکرار رفتار ارتباط دارد.

برای مثال اگر فرد بعد از کار چای بنوشد و این کار برای او تجربهٔ خوشایندی باشد، سیستم \eng{Mesolimbic} می‌تواند به این الگو شرطی شود. در نتیجه، بعد از کار میل به نوشیدن چای افزایش پیدا می‌کند. به زبان ساده، این مسیر به مغز می‌گوید: این رفتار قبلاً خوب بوده است؛ دوباره انجامش بده.

در توضیح این بخش، مسیر \eng{Mesolimbic} از \eng{VTA} به نواحی لیمبیک، به‌ویژه \eng{Nucleus Accumbens}، می‌رسد. این مسیر بیشتر با تقویت یا تضعیف ارزش محرک‌ها در حافظهٔ بلندمدت ارتباط دارد. وقتی خطای پیش‌بینی پاداش مثبت است، این مسیر کمک می‌کند محرک یا عمل مربوطه در حافظه با ارزش بیشتری ثبت شود.

\subsubsection{\eng{Mesocortical Pathway}}
مسیر \eng{Mesocortical} نیز از \eng{VTA} شروع می‌شود، اما به سمت \eng{Prefrontal Cortex} یا قشر پیش‌پیشانی می‌رود. قشر پیش‌پیشانی با تصمیم‌گیری، تحلیل، ارزیابی پیامدها، کنترل رفتار، بررسی ارزش تلاش و فکر کردن به نتایج بلندمدت ارتباط دارد.

مثلاً وقتی فرد کیک می‌بیند، مسیر \eng{Mesolimbic} ممکن است او را به سمت خوردن کیک سوق دهد، چون کیک قبلاً تجربه‌ای لذت‌بخش بوده است. اما مسیر \eng{Mesocortical} پیامدها را ارزیابی می‌کند: آیا خوردن کیک منطقی است؟ آیا فرد گرسنه است؟ آیا رژیم دارد؟ آیا نتیجهٔ بلندمدت این رفتار خوب است؟ در این بیان استعاری، \eng{Mesolimbic} بیشتر نقش گاز و \eng{Mesocortical} بیشتر نقش تحلیل‌گر یا ترمز را دارد.

\subsection{اختلال‌های مرتبط با سیستم پاداش}
در ادامه، چند پدیده با همین چارچوب توضیح داده شد. هدف این بخش تشخیص پزشکی نیست؛ بلکه نشان دادن این است که نظریهٔ پاداش، دوپامین و خطای پیش‌بینی چگونه می‌تواند برای فهم بعضی تغییرات رفتاری به‌کار رود.

\subsubsection{پیری}
در پیری، سیستم پاداش معمولاً کم‌تر دچار غافلگیری می‌شود. فرد تجربه‌های زیادی کسب کرده و بسیاری از محرک‌ها برای او قابل پیش‌بینی شده‌اند. بنابراین خطای پیش‌بینی پاداش اغلب به صفر نزدیک‌تر است، دوپامین کم‌تر ترشح می‌شود، به‌روزرسانی مدل ذهنی کم‌تر رخ می‌دهد و میل به تجربهٔ چیزهای تازه یا \eng{exploration} کاهش می‌یابد.

از این دید، یکی از راه‌های ایجاد حال بهتر برای فرد سالمند می‌تواند فراهم کردن تجربه‌های تازه و کمتر تجربه‌شده باشد؛ مانند سفرهای جدید، فعالیت‌های جدید یا موقعیت‌هایی که پیش‌تر در زندگی فرد وجود نداشته‌اند.

\subsubsection{افسردگی}
افسردگی نیز در این بحث با همین زبان توضیح داده شد. در این حالت، سیستم پاداش به‌درستی فعال نمی‌شود و فرد ممکن است تجربه کند که چیزی خوشحالش نمی‌کند، کارهایی که قبلاً لذت‌بخش بودند دیگر جذاب نیستند، یا انجام دادن هیچ کاری ارزش ندارد.

در این چارچوب، فرد ممکن است در حالت قطعیت بماند و عمل تازه‌ای انجام ندهد. وقتی تجربهٔ تازه‌ای رخ ندهد، خطای پیش‌بینی پاداش نیز صفر می‌ماند. اما این با یادگیری سالم فرق دارد. در یادگیری سالم، فرد محیط را تجربه می‌کند تا مدل ذهنی‌اش دقیق‌تر شود و خطا کاهش یابد؛ اما در این توصیف از افسردگی، فرد با انجام ندادن عمل و تجربه نکردن، به‌صورت ناسالم به خطای نزدیک به صفر می‌رسد.

\subsubsection{اعتیاد}
اعتیاد یکی از مثال‌های اصلی سوءاستفاده از سیستم پاداش است. در ابتدای مصرف، محرک اعتیادآور می‌تواند خطای پیش‌بینی پاداش مثبت ایجاد کند:
\[
\text{\lr{RPE}} > 0
\]
چون تجربه قوی‌تر یا متفاوت‌تر از انتظار قبلی است. اما پس از تکرار، سیستم یاد می‌گیرد و آن تجربه برای مغز عادی‌تر می‌شود؛ بنابراین خطای پیش‌بینی پاداش به سمت صفر می‌رود:
\[
\text{\lr{RPE}} \rightarrow 0
\]

در اینجا چرخهٔ معیوب شروع می‌شود. فرد برای اینکه دوباره خطای مثبت و پاسخ دوپامینی ایجاد کند، دوز را بالا می‌برد. با افزایش دوز دوباره غافلگیری ایجاد می‌شود، اما پس از مدتی مغز دوباره سازگار می‌شود و خطا کاهش می‌یابد. چرخهٔ کلی چنین است: تجربهٔ اولیه خطای مثبت ایجاد می‌کند؛ سیستم یاد می‌گیرد؛ خطا به صفر نزدیک می‌شود؛ دوز افزایش می‌یابد؛ دوباره خطای مثبت ایجاد می‌شود؛ و مغز دوباره سازگار می‌شود.

مغز در برابر این چرخه منفعل نمی‌ماند. یکی از واکنش‌های دفاعی آن کاهش حساسیت به دوپامین است. در نتیجه، فرد نه‌تنها از محرک اعتیادآور کم‌تر لذت می‌برد، بلکه ممکن است از محرک‌های عادی زندگی نیز لذت کم‌تری تجربه کند.

در این بحث برای این پرسش که چرا همهٔ افراد وارد چرخهٔ اعتیاد نمی‌شوند، چند عامل مطرح شد. برخی افراد اصلاً در معرض محرکی قرار نمی‌گیرند که خطای مثبت شدید ایجاد کند. آگاهی و محیط نیز مهم‌اند، چون شناخت پیامدهای اعتیاد می‌تواند از ورود به چرخه جلوگیری کند. عامل دیگر تفاوت‌های فردی، ژنتیک و قدرت مسیر \eng{Mesocortical} است. اگر این مسیر در کنترل، راهبرد و ارزیابی پیامدهای بلندمدت قوی‌تر عمل کند، فرد بهتر می‌تواند در برابر پاداش فوری مقاومت کند.

مثال مطرح‌شده \eng{Marshmallow Test} بود: آزمایشی که در آن به کودک گفته می‌شود اگر اکنون شیرینی یا شکلات را نخورد، بعداً پاداش بیشتری می‌گیرد. برخی کودکان می‌توانند صبر کنند و برخی نه. این مثال برای توضیح خودکنترلی و توانایی دیدن پیامدهای بلندمدت استفاده شد.

\subsection{ارتباط سیستم پاداش با \eng{Predictive Coding}}
در این بخش پرسیده شد کدام نظریهٔ شناختی می‌تواند این فرایند را توضیح دهد. پاسخ اصلی \eng{Predictive Coding} بود. در کدگذاری پیش‌بینانه، سیستم همواره میان دانش قبلی یا \eng{prior knowledge} و واقعیت مشاهده‌شده مقایسه انجام می‌دهد. آنچه برای یادگیری مهم است، صرفاً خود واقعیت خام نیست؛ بلکه خطای میان پیش‌بینی و واقعیت است.

در سیستم پاداش نیز همین اتفاق می‌افتد. عامل پیش از عمل، انتظاری از پاداش دارد. سپس عمل انجام می‌دهد و پاداش واقعی را دریافت می‌کند. اختلاف میان پاداش واقعی و پاداش پیش‌بینی‌شده همان چیزی است که باید یاد گرفته شود.

\subsection{\eng{Reward Prediction Error}}
مفهوم مرکزی این بخش \eng{Reward Prediction Error} یا خطای پیش‌بینی پاداش است:
\[
\text{\lr{Reward Prediction Error}} =
\text{\lr{Actual Reward}} - \text{\lr{Predicted Reward}}
\]
یا به‌شکل خلاصه:
\[
\text{\lr{RPE}} = R_{\text{\lr{actual}}} - R_{\text{\lr{predicted}}}
\]

اگر پاداش واقعی بیشتر از پاداش پیش‌بینی‌شده باشد، خطا مثبت است:
\[
\text{\lr{RPE}} > 0
\]
یعنی نتیجه بهتر از انتظار بوده است. در این حالت سیستم غافلگیر می‌شود و عمل یا محرک مربوطه تقویت می‌شود. مثال رستورانی بود که فرد انتظار متوسطی از آن دارد، اما تجربهٔ واقعی بسیار بهتر از انتظار اوست. در این حالت، ارزش آن رستوران در ذهن افزایش پیدا می‌کند.

اگر پاداش واقعی تقریباً برابر با پیش‌بینی قبلی باشد، خطا نزدیک به صفر است:
\[
\text{\lr{RPE}} = 0
\]
در این حالت اتفاق تازه‌ای برای یادگیری رخ نمی‌دهد و مدل قبلی تقریباً ثابت می‌ماند.

اگر پاداش واقعی کم‌تر از پاداش پیش‌بینی‌شده باشد، خطا منفی است:
\[
\text{\lr{RPE}} < 0
\]
یعنی نتیجه بدتر از انتظار بوده است. در این حالت سیستم باید یاد بگیرد که ارزش آن عمل یا محرک را کاهش دهد.

\subsection{\eng{Free Energy Principle} و کاهش \eng{Surprise}}
در ادامه، \eng{Free Energy Principle} مطرح شد. بر اساس این نگاه، انسان فقط به‌دنبال دریافت پاداش بیشتر نیست؛ بلکه تلاش می‌کند عدم قطعیت و خطای پیش‌بینی را کاهش دهد. سیستم شناختی محیط را تجربه می‌کند تا مدل ذهنی‌اش دقیق‌تر شود و اختلاف میان پیش‌بینی و واقعیت کم‌تر شود.

هدف نهایی این است که خطای پیش‌بینی پاداش به صفر نزدیک شود:
\[
\text{\lr{RPE}} \rightarrow 0
\]
یعنی سیستم به حالتی برسد که پاداش واقعی با پاداش پیش‌بینی‌شده همخوان شود. در این حالت مدل ذهنی به‌روزرسانی شده و غافلگیری زیادی رخ نمی‌دهد.

\section{مدل‌های محاسباتی یادگیری تقویتی}
از دید محاسباتی، سیستم پاداش انسان شباهت زیادی به \eng{Reinforcement Learning} دارد. در یادگیری تقویتی، عامل در محیط قرار دارد، عملی انجام می‌دهد، از یک وضعیت به وضعیت دیگر می‌رود، پاداش دریافت می‌کند و بر اساس آن پاداش یاد می‌گیرد در آینده چه رفتاری انجام دهد.

\subsection{مسئلهٔ \eng{Reinforcement Learning}}
مسئلهٔ اصلی \eng{Reinforcement Learning} یادگیری از طریق تعامل با محیط است. عامل برخلاف مدل نظارت‌شده، از ابتدا برچسب درست هر تصمیم را ندارد؛ باید عملی انجام دهد و پیامد آن را ببیند. اگر پاداش دریافت‌شده مطلوب باشد، احتمال تکرار آن رفتار بیشتر می‌شود و اگر پیامد نامطلوب باشد، احتمال آن رفتار کاهش می‌یابد.

یادگیری تقویتی وقتی اهمیت ویژه پیدا می‌کند که محیط آنلاین باشد؛ یعنی عامل مدل آماده و دقیقی از محیط ندارد و باید تجربه را نمونه‌به‌نمونه به‌دست آورد. در هر گام، عامل یک \eng{state} را مشاهده می‌کند، یک \eng{action} انجام می‌دهد، به وضعیت جدیدی منتقل می‌شود و از محیط \eng{reward} یا \eng{punishment} دریافت می‌کند. سپس مدل خود را بر اساس همین تجربه به‌روزرسانی می‌کند.

حالت مهم دیگر، محیط تصادفی یا \eng{stochastic} است. در چنین محیطی، اگر در یک وضعیت مشخص یک عمل مشخص انجام شود، لزوماً همیشه به یک وضعیت مشخص نمی‌رسیم. گذار کلی چنین است:
\[
S \xrightarrow{A} S'
\]
اما \(S'\) می‌تواند چند حالت مختلف داشته باشد؛ برای مثال ممکن است با احتمال ۰٫۵ به \(S'_1\) و با احتمال ۰٫۲ به \(S'_2\) برسیم. چون مدل گذار از ابتدا معلوم نیست، عامل باید با تجربه کردن و گرفتن بازخورد آن را یاد بگیرد.

\subsection{تفاوت با یادگیری نظارت‌شده و بدون‌ناظر}
در \eng{Supervised Learning}، داده‌ها برچسب دارند و مدل از رابطهٔ میان \eng{input} و \eng{label} یاد می‌گیرد. اما در سیستم پاداش و یادگیری تقویتی، برچسب آماده وجود ندارد. عامل باید \eng{action} انجام دهد، از محیط \eng{feedback} بگیرد و بر اساس نتیجهٔ تجربه یاد بگیرد.

در \eng{Unsupervised Learning} نیز برچسب مستقیم وجود ندارد، اما مدل معمولاً از ساختار داده‌ها الگو استخراج می‌کند؛ مثلاً خوشه‌بندی می‌کند یا بازنمایی می‌آموزد. تفاوت \eng{Reinforcement Learning} این است که مسئله فقط نبودن برچسب نیست. نکتهٔ اصلی تعامل فعال با محیط و دریافت \eng{reward} پس از عمل است.

\subsection{\eng{Temporal Difference Learning}}
در پایان این بخش، بحث به فرمول یادگیری رسید. فرض کنیم عامل در وضعیت فعلی \(S\) قرار دارد و با انجام عمل \(A\) به وضعیت بعدی \(S'\) می‌رود. عامل برای وضعیت فعلی یک ارزش قبلی دارد که آن را \(V_t\) می‌نامیم. پس از تجربهٔ جدید، یک نمونهٔ تازه از ارزش آن تجربه به‌دست می‌آید. اکنون سیستم باید ارزش قبلی را با توجه به تجربهٔ جدید به‌روزرسانی کند.

نه دانش قبلی به‌تنهایی کافی است و نه نمونهٔ جدید به‌تنهایی. بنابراین سیستم میان این دو درون‌یابی می‌کند. برای این کار از ضریب \(\alpha\) استفاده می‌شود:
\[
0 \leq \alpha \leq 1
\]
\(\alpha\) همان \eng{learning rate} است. اگر \(\alpha = 1\) باشد، فقط تجربهٔ جدید مهم است و دانش قبلی نادیده گرفته می‌شود. اگر \(\alpha = 0\) باشد، تجربهٔ جدید اثری ندارد و سیستم چیزی یاد نمی‌گیرد.

فرمول کلی به‌روزرسانی چنین است:
\[
V_{t+1} = V_t + \alpha(\text{\lr{Sample}} - V_t)
\]
ترم داخل پرانتز همان خطای پیش‌بینی پاداش است:
\[
\text{\lr{Sample}} - V_t = \text{\lr{RPE}}
\]
پس:
\[
V_{t+1} = V_t + \alpha \times \text{\lr{RPE}}
\]

این منطق \eng{Temporal Difference Learning} نامیده می‌شود، چون اختلاف میان دو زمان را محاسبه می‌کند: پیش از عمل در زمان \(t\) و پس از عمل در زمان \(t+1\). سیستم ارزش قبلی را به اندازهٔ نرخ یادگیری و بر اساس خطای پیش‌بینی به‌روزرسانی می‌کند.

هدف نهایی این فرایند آن است که نمونهٔ تجربه‌شده با ارزش پیش‌بینی‌شده همگرا شود:
\[
\text{\lr{Sample}} = V_t
\]
در نتیجه:
\[
\text{\lr{RPE}} = 0
\]
این همان ایده‌ای است که در بحث \eng{Free Energy Principle} نیز دیده شد: کاهش غافلگیری و کاهش خطای پیش‌بینی.

\subsection{\eng{Discount Factor}}
در این بخش، هدف اصلی عمیق‌تر کردن مفاهیم بخش‌های قبل بود؛ به‌ویژه \eng{Q-Learning}، \eng{Active RL} و ارتباط آن‌ها با سیستم پاداش. نقطهٔ شروع این بود که عامل در یک فضای جست‌وجو حرکت می‌کند. ممکن است دربارهٔ بعضی حالت‌ها تجربهٔ قبلی یا \eng{prior knowledge} داشته باشد، و دربارهٔ بعضی حالت‌ها هیچ تجربه‌ای نداشته باشد؛ در این حالت، مقدار اولیه معمولاً صفر در نظر گرفته می‌شود.

فرض کنید عامل در یک حالت شروع قرار دارد و دو انتخاب دارد: رفتن به چپ یا رفتن به راست. اگر عامل فقط پاداش فوری را ببیند، ممکن است بگوید چپ پاداش ۱ دارد و راست فعلاً پاداش ۰ دارد؛ پس در حالت کاملاً \eng{greedy} به چپ می‌رود. اما ممکن است رفتن به راست در آینده به پاداش بزرگ‌تری، مثلاً ۱۰، برسد. بنابراین عامل فقط نباید به پاداش فوری نگاه کند؛ باید تا حدی آینده را نیز در نظر بگیرد.

برای مدل کردن میزان اهمیت آینده از \eng{Discount Factor} یا ضریب تخفیف استفاده می‌کنیم. این ضریب معمولاً با \(\gamma\) نشان داده می‌شود:
\[
0 \leq \gamma \leq 1
\]
اگر \(\gamma = 0\) باشد، عامل فقط پاداش فوری را می‌بیند و کاملاً کوتاه‌نگر است. اگر \(\gamma = 1\) باشد، عامل آینده را بدون تخفیف کامل لحاظ می‌کند. اگر \(0 < \gamma < 1\) باشد، پاداش‌های آینده در محاسبه وارد می‌شوند، اما هرچه دورتر باشند اثر کم‌تری دارند:
\[
r_t + \gamma r_{t+1} + \gamma^2 r_{t+2} + \gamma^3 r_{t+3} + \cdots
\]

نکتهٔ مهم دیگر این است که مفهوم ضریب تخفیف را می‌توان تا حدی با سازوکارهای زیستی نیز مرتبط دانست. برای مثال دربارهٔ سروتونین گفته شد که اگر سطح آن بالاتر باشد، فرد ممکن است کم‌تر عجولانه و کم‌تر کاملاً \eng{greedy} تصمیم بگیرد و توان بیشتری برای در نظر گرفتن آینده داشته باشد. البته این رفتار فقط به یک مادهٔ زیستی وابسته نیست و ترکیبی از عوامل مختلف می‌تواند روی چیزی شبیه ضریب تخفیف اثر بگذارد.

\subsection{\eng{Value Function}}
در \eng{TD Learning} هدف به‌روزرسانی مقدار یک حالت است:
\[
V(s)
\]
فرمول کامل‌تر به‌روزرسانی با ضریب تخفیف چنین است:
\[
V_{t+1}(s) = (1-\alpha)V_t(s) + \alpha \left[r + \gamma V_t(s')\right]
\]
در این فرمول، \(V_t(s)\) مقدار قبلی حالت \(s\)، \(r\) پاداش دریافت‌شده، \(s'\) حالت بعدی، \(V_t(s')\) مقدار حالت بعدی، \(\alpha\) نرخ یادگیری و \(\gamma\) ضریب تخفیف است.

این فرمول نوعی درون‌یابی میان دانش قبلی و تجربهٔ جدید است. بخش \((1-\alpha)V_t(s)\) دانش قبلی را نگه می‌دارد و بخش \(\alpha[r+\gamma V_t(s')]\) تجربهٔ جدید را وارد می‌کند.

مثال سه تجربه داشت:
\[
A \xrightarrow{a_1,\, r=-3} B
\]
\[
B \xrightarrow{a_1,\, r=4} A
\]
\[
A \xrightarrow{a_2,\, r=-4} A
\]
فرض شد \(\alpha=0.1\)، \(\gamma=1\)، و در ابتدا \(V(A)=0\) و \(V(B)=0\).

در تجربهٔ اول، از \(A\) به \(B\) با پاداش \(-3\) می‌رویم:
\[
V(A)=0.9V(A)+0.1[-3+V(B)]
\]
چون مقدارهای اولیه صفرند:
\[
V(A)=0.9(0)+0.1[-3+0]=-0.3
\]
پس مقدار \(A\) کاهش پیدا می‌کند. در تجربهٔ دوم:
\[
V(B)=0.9V(B)+0.1[4+V(A)]
\]
و چون \(V(A)=-0.3\):
\[
V(B)=0.9(0)+0.1[4-0.3]=0.37
\]
در تجربهٔ سوم دوباره \(V(A)\) به‌روزرسانی می‌شود:
\[
V(A)=0.9(-0.3)+0.1[-4-0.3]
\]
نکتهٔ اصلی این است که در \eng{TD Learning}، در هر مرحله مقدار همان حالتی به‌روزرسانی می‌شود که عامل از آن حرکت کرده است.

\subsection{\eng{Q Function}}
تا اینجا مقدار حالت‌ها را داشتیم، اما هدف نهایی عامل فقط دانستن ارزش حالت‌ها نیست. عامل باید در هر لحظه تصمیم بگیرد کدام عمل را انتخاب کند. بنابراین لازم است بدانیم اگر عامل در حالت \(s\) باشد و عمل \(a\) را انجام دهد، این انتخاب چقدر ارزش دارد.

این مقدار با \(Q(s,a)\) نشان داده می‌شود. \(Q(s,a)\) ارزش انجام عمل \(a\) در حالت \(s\) است. رابطهٔ کلی میان \(V\) و \(Q\) چنین است:
\[
V(s)=\max_a Q(s,a)
\]
یعنی ارزش یک حالت برابر است با بیشترین مقدار \(Q\) میان عمل‌های ممکن در آن حالت.

در دنیای واقعی، اگر یک عمل را انتخاب کنیم، همیشه دقیقاً به همان حالتی که انتظار داریم نمی‌رسیم. برای مثال ممکن است عامل تصمیم بگیرد به چپ برود؛ با احتمال ۰٫۹ واقعاً به خانهٔ چپ برسد، اما با احتمال ۰٫۱ به حالت دیگری منتقل شود. این احتمال گذار با تابع \(T(s,a,s')\) نشان داده می‌شود:
\[
T(s,a,s')
\]
یعنی احتمال اینکه عامل از حالت \(s\)، با انجام عمل \(a\)، به حالت \(s'\) برود. در این حالت می‌توان مقدار \(Q\) را با جمع روی حالت‌های بعدی ممکن نوشت:
\[
Q(s,a)=\sum_{s'} T(s,a,s')\left[R(s,a,s')+\gamma V(s')\right]
\]
در این رابطه، عامل هم پاداش فوری را در نظر می‌گیرد و هم ارزش حالت بعدی را، اما ارزش آینده با ضریب تخفیف وارد می‌شود.

\subsection{\eng{Q-Learning}}
در \eng{Q-Learning} به‌جای اینکه فقط \(V(s)\) به‌روزرسانی شود، مستقیماً \(Q(s,a)\) به‌روزرسانی می‌شود:
\[
Q_{t+1}(s,a) = (1-\alpha)Q_t(s,a)
+ \alpha \left[r+\gamma \max_{a'}Q_t(s',a')\right]
\]
در این فرمول، \(s\) حالت فعلی، \(a\) عمل انجام‌شده، \(r\) پاداش دریافت‌شده، \(s'\) حالت بعدی، و \(a'\) عمل‌های ممکن در حالت بعدی هستند. عبارت \(\max_{a'}Q_t(s',a')\) بهترین عمل ممکن در حالت بعدی را نشان می‌دهد.

تفاوت مهم \eng{Q-Learning} با \eng{TD Learning} ساده این است که در \eng{Q-Learning} مستقیماً ارزش عمل‌ها در حالت‌ها یاد گرفته می‌شود، نه فقط ارزش خود حالت‌ها. به همین دلیل، \eng{Q-Learning} یک گام به تصمیم‌گیری واقعی نزدیک‌تر است.

برای مثال، همان اپیزود قبل را در نظر می‌گیریم:
\[
A \xrightarrow{a_1,\, r=-3} B,\qquad
B \xrightarrow{a_1,\, r=4} A,\qquad
A \xrightarrow{a_2,\, r=-4} A
\]
فرض می‌کنیم \(\alpha=0.1\)، \(\gamma=1\)، و همهٔ مقدارهای \(Q\) در ابتدا صفرند.

در تجربهٔ اول، عامل از \(A\) با عمل ۱ به \(B\) می‌رود و پاداش \(-3\) می‌گیرد:
\[
Q(A,1)=0.9Q(A,1)+0.1[-3+\max_{a'}Q(B,a')]
\]
چون هنوز دربارهٔ عمل‌های حالت \(B\) چیزی نمی‌دانیم، بیشینه صفر است:
\[
Q(A,1)=0.1[-3]=-0.3
\]

در تجربهٔ دوم، عامل از \(B\) با عمل ۱ به \(A\) می‌رود و پاداش ۴ می‌گیرد:
\[
Q(B,1)=0.9Q(B,1)+0.1[4+\max_{a'}Q(A,a')]
\]
در حالت \(A\)، مقدار \(Q(A,1)\) برابر \(-0.3\) است و عمل‌های دیگر هنوز مقدار صفر دارند؛ پس بیشینه صفر می‌شود:
\[
Q(B,1)=0.1[4+0]=0.4
\]

در تجربهٔ سوم، عامل از \(A\) با عمل ۲ به \(A\) می‌رود و پاداش \(-4\) می‌گیرد:
\[
Q(A,2)=0.9Q(A,2)+0.1[-4+\max_{a'}Q(A,a')]
\]
در این مرحله، بیشینهٔ مقدارهای \(Q\) در \(A\) همچنان صفر است، پس:
\[
Q(A,2)=0.1[-4+0]=-0.4
\]
روند به همین شکل ادامه پیدا می‌کند. در هر تجربه، مقدار \(Q(s,a)\) مربوط به همان حالت و همان عمل به‌روزرسانی می‌شود.

\subsection{محیط‌های تصادفی و عدم قطعیت}
در محیط قطعی، انجام یک عمل مشخص در یک وضعیت مشخص همیشه به نتیجه‌ای مشخص می‌رسد. اما در محیط تصادفی یا غیرقطعی، یک عمل می‌تواند نتایج مختلفی داشته باشد. مثال رانندگی بود: در یک محیط ایده‌آل، اگر راننده فرمان را به چپ بچرخاند، انتظار داریم ماشین به چپ برود؛ اما در محیطی غیرقطعی ممکن است ماشین کمی منحرف شود، متوقف شود یا رفتار دیگری نشان دهد.

بنابراین عامل باید با عدم قطعیت در گذارها کنار بیاید. یادگیری در این حالت یعنی تجربه کردن، دریافت بازخورد و اصلاح تدریجی مدل از اینکه هر عمل در هر وضعیت معمولاً به چه پیامدهایی منجر می‌شود.

\subsection{\eng{Exploration} و \eng{Exploitation}}
اگر عامل فقط از رابطهٔ
\[
V(s)=\max_a Q(s,a)
\]
برای انتخاب عمل استفاده کند، همیشه عملی را انتخاب می‌کند که تا آن لحظه بهترین تجربهٔ شناخته‌شده را داشته است. این حالت کاملاً \eng{greedy} یا \eng{exploitative} است؛ یعنی عامل فقط از تجربه‌های قبلی خود استفاده می‌کند.

اما در تصمیم‌گیری واقعی، همیشه فقط بهترین تجربهٔ قبلی را تکرار نمی‌کنیم. مثال خرید این نکته را روشن می‌کند: معمولاً همیشه فقط از جایی خرید نمی‌کنیم که قبلاً بهترین تجربه را از آن داشته‌ایم؛ گاهی مکان‌های جدید را هم امتحان می‌کنیم، به‌ویژه اگر فکر کنیم شاید گزینهٔ بهتری وجود داشته باشد.

\eng{Exploitation} یعنی عامل از بهترین تجربه‌های قبلی خود استفاده کند: «تا الان این عمل بهترین بوده، پس دوباره همان را انتخاب می‌کنم.» در مقابل، \eng{Exploration} یعنی عامل به عمل‌هایی هم فرصت بدهد که کم‌تر امتحان شده‌اند: «این عمل هنوز بهترین شناخته نشده، اما چون کم‌تر امتحان شده، شاید ارزش کشف کردن داشته باشد.» تصمیم‌گیری در دنیای واقعی معمولاً ترکیبی از این دو است.

\subsection{\eng{Passive RL} و \eng{Active RL}}
برای اینکه عامل فقط \eng{greedy} نباشد، باید عدم قطعیت را هم وارد تصمیم‌گیری کنیم. اگر امتیاز یک عمل فقط برابر با مقدار \(Q\) باشد:
\[
\text{\lr{Score}}(s,a)=Q(s,a)
\]
عامل بیشتر به سمت بهره‌برداری از تجربه‌های قبلی می‌رود. اما می‌توان جمله‌ای مربوط به عدم قطعیت نیز اضافه کرد:
\[
\text{\lr{Score}}(s,a)=Q(s,a)+\beta U(s,a)
\]
در این رابطه، \(Q(s,a)\) تجربهٔ قبلی و پاداش یادگرفته‌شده را نشان می‌دهد، \(U(s,a)\) میزان عدم قطعیت دربارهٔ آن عمل است، و \(\beta\) مشخص می‌کند عامل چقدر به عدم قطعیت اهمیت می‌دهد.

اگر عملی کم‌تر امتحان شده باشد، عدم قطعیت آن بیشتر است. بنابراین حتی اگر مقدار \(Q\) آن فعلاً خیلی بالا نباشد، ممکن است به خاطر عدم قطعیت بالا انتخاب شود تا عامل آن را کشف کند. این ایده به روش‌هایی مانند \eng{Upper Confidence Bound} یا \eng{UCB} شبیه است؛ یعنی عامل فقط میانگین تجربه‌شده را نگاه نمی‌کند، بلکه برای گزینه‌های کم‌تر امتحان‌شده یک حد بالای خوش‌بینانه هم در نظر می‌گیرد.

در حالت \eng{passive}، عامل بیشتر بر اساس تجربه‌های قبلی تصمیم می‌گیرد؛ بهترین گزینه‌های شناخته‌شده را انتخاب می‌کند، تمایل کمی به امتحان گزینه‌های جدید دارد و بیشتر \eng{exploit} می‌کند. در حالت \eng{active}، عامل علاوه بر تجربه‌های قبلی، عدم قطعیت را هم در نظر می‌گیرد؛ گزینه‌های کم‌تر امتحان‌شده را بررسی می‌کند و میان \eng{exploration} و \eng{exploitation} تعادل برقرار می‌کند.

برای توضیح شهودی گفت می‌توان حالت \eng{active} را شبیه جوانی دانست که بیشتر دوست دارد تجربه و کشف کند، و حالت \eng{passive} را شبیه پیری دانست که بیشتر به تجربه‌های قبلی اتکا دارد. این تشبیه برای فهم شهودی است، نه یک ادعای زیستی دقیق.
